\section{Introducción}
Las ondas estacionarias acústicas permiten relacionar las frecuencias naturales de una columna de aire con su geometría y con la velocidad del sonido. La superposición de ondas incidentes y reflejadas genera modos compatibles con las condiciones de frontera; la resonancia facilita su detección mediante máximos en un espectro de frecuencias \cite{halliday2018,openstax2016}.

Una probeta constituye un resonador cerrado por su base y abierto por su boca. Al incorporar agua, la superficie del líquido establece el cierre inferior de la columna de aire y permite variar su longitud. Esta configuración ofrece dos maneras de estimar la velocidad del sonido: relacionar las frecuencias con los armónicos permitidos y estudiar la frecuencia fundamental frente a la longitud inversa.

De acuerdo con la guía oficial \cite{guiauca}, los objetivos son generar y detectar ondas estacionarias e identificar sus resonancias; analizar su relación con los armónicos permitidos; determinar la velocidad del sonido mediante la frecuencia fundamental y la longitud efectiva, comparándola con la referencia térmica; y estudiar la relación entre longitud de aire y frecuencia de resonancia. El análisis aborda estos objetivos mediante las resonancias registradas e incluye las predicciones de volumen y las observaciones cualitativas obtenidas con cuatro diapasones.
